\documentclass[11pt,a4paper]{article}
\usepackage[utf8]{inputenc}
\usepackage[T1]{fontenc}
\usepackage[brazil]{babel}
\usepackage{amsmath,amssymb}
\usepackage{graphicx}
\usepackage{geometry}
\usepackage{enumitem}
\usepackage{array,booktabs,tabularx}
\usepackage{xcolor}
\usepackage{fancyhdr}
\usepackage{microtype}
\usepackage[most]{tcolorbox}
\usepackage{helvet}
\usepackage{caption}
\usepackage{float}
\usepackage{hyperref}
\renewcommand{\familydefault}{\sfdefault}

\geometry{top=20mm,bottom=20mm,left=22mm,right=22mm}
\definecolor{unirblue}{HTML}{0B4EA2}
\definecolor{ink}{HTML}{172033}
\definecolor{muted}{HTML}{667085}
\definecolor{soft}{HTML}{F5F7FA}
\definecolor{line}{HTML}{D7DCE2}
\definecolor{guide}{HTML}{EEF4FB}

\hypersetup{colorlinks=true,linkcolor=unirblue,urlcolor=unirblue,citecolor=unirblue}
\setlength{\parindent}{0pt}
\setlength{\parskip}{5pt}
\setlist[itemize]{leftmargin=5mm,itemsep=2pt,topsep=2pt}
\setlist[enumerate]{leftmargin=6mm,itemsep=2pt,topsep=2pt}
\captionsetup{font=small,labelfont=bf}

\pagestyle{fancy}
\fancyhf{}
\renewcommand{\headrulewidth}{0pt}
\cfoot{\color{muted}\footnotesize UNIR · Processamento de Imagens · Relatório da Avaliação \hfill \thepage}

\newcommand{\sectiontitle}[1]{%
  \vspace{5pt}{\Large\bfseries\color{ink}#1}\par
  \vspace{1mm}{\color{unirblue}\rule{\textwidth}{0.8pt}}\vspace{2mm}}
\newcommand{\subsectiontitle}[1]{\vspace{4pt}{\large\bfseries\color{ink}#1}\par\vspace{1mm}}
\newcommand{\placeholder}[1]{\textcolor{muted}{\textit{#1}}}

\newtcolorbox{orientacao}{
  colback=guide,
  colframe=unirblue!25,
  leftrule=2pt,rightrule=0pt,toprule=0pt,bottomrule=0pt,
  arc=0mm,boxsep=1.5mm,left=2mm,right=2mm,top=1.5mm,bottom=1.5mm,
  fontupper=\small\color{ink}
}

\newtcolorbox{campo}{
  colback=soft,colframe=line,boxrule=.4pt,arc=1mm,
  left=2mm,right=2mm,top=2mm,bottom=2mm
}

\begin{document}

% ============================================================
% CAPA
% ============================================================
\begin{titlepage}
\thispagestyle{empty}

\begin{center}
  \includegraphics[width=4.8cm]{unir_logo.png}
\end{center}

\vspace{8mm}

\begin{center}
{\small\bfseries\color{muted} FUNDAÇÃO UNIVERSIDADE FEDERAL DE RONDÔNIA - UNIR}\\[2mm]
{\small\color{muted} Faculdade de Engenharia, Tecnologia e Inovação - FAETI}\\
{\small\color{muted} Bacharelado em Ciência da Computação}
\end{center}

\vfill

\begin{center}
{\small\bfseries\color{unirblue} PROCESSAMENTO DE IMAGENS}\\[5mm]
{\Huge\bfseries\color{ink} Relatório da Primeira Avaliação}\\[3mm]
{\large\color{muted} Sistema de Processamento de Imagens}\\[12mm]

{\Large\bfseries\color{ink} NOME COMPLETO DO ALUNO}\\[2mm]
{\small\color{muted} Matrícula: 00000000000}
\end{center}

\vfill

\begin{center}
{\small\color{muted} Professor: Lucas Marques da Cunha}\\[1mm]
{\small\color{muted} Porto Velho - RO}\\
{\small\color{muted} 2026}
\end{center}

\end{titlepage}

% ============================================================
% FOLHA DE IDENTIFICAÇÃO
% ============================================================
\sectiontitle{Identificação do trabalho}

\begin{campo}
\begin{tabularx}{\textwidth}{>{\bfseries}p{37mm}X}
Aluno(a): & \placeholder{Nome completo}\\[2mm]
Matrícula: & \placeholder{Número de matrícula}\\[2mm]
Disciplina: & Processamento de Imagens\\[2mm]
Professor: & Lucas Marques da Cunha\\[2mm]
Título do sistema: & \placeholder{Nome atribuído ao sistema desenvolvido}\\[2mm]
Linguagem/ambiente: & \placeholder{Ex.: Python/desktop, Java/web, C++/desktop etc.}\\[2mm]
Repositório: & \placeholder{URL do repositório GitHub}\\
\end{tabularx}
\end{campo}

\begin{orientacao}
\textbf{Como usar este template.} Substitua todos os textos em itálico/cinza pelas informações do seu trabalho. As caixas azuis contêm orientações e devem ser \textbf{removidas da versão final}. O relatório deve documentar o que foi efetivamente implementado e testado no sistema entregue.
\end{orientacao}

\vspace{3mm}
\tableofcontents
\newpage

% ============================================================
% 1 INTRODUÇÃO
% ============================================================
\section{Introdução}

\begin{orientacao}
Apresente o contexto do trabalho e o problema abordado. Explique, de forma breve, o que é processamento digital de imagens e por que operações como transformações de intensidade, histogramas, filtragem, convolução, detecção de bordas e aguçamento são relevantes. Finalize delimitando o sistema desenvolvido. Evite transformar a introdução em uma lista de funcionalidades.
\end{orientacao}

\placeholder{Escreva aqui a introdução do trabalho. Recomenda-se contextualizar o tema em 3 a 5 parágrafos.}

\subsection{Objetivos}

\begin{orientacao}
Apresente um objetivo geral e objetivos específicos coerentes com a avaliação. Os objetivos específicos devem refletir as principais funcionalidades implementadas e avaliadas.
\end{orientacao}

\textbf{Objetivo geral.} \placeholder{Descreva o objetivo principal do sistema.}

\textbf{Objetivos específicos.}
\begin{itemize}
  \item \placeholder{Objetivo específico 1.}
  \item \placeholder{Objetivo específico 2.}
  \item \placeholder{Objetivo específico 3.}
\end{itemize}

% ============================================================
% 2 FUNDAMENTAÇÃO
% ============================================================
\section{Fundamentação teórica}

\begin{orientacao}
Explique os conceitos necessários para compreender os algoritmos implementados. Não copie definições isoladas: apresente as equações, descreva o significado dos parâmetros e relacione cada conceito ao comportamento esperado na imagem. Cite as referências utilizadas.
\end{orientacao}

\subsection{Representação de imagens digitais e vizinhança}
\placeholder{Explique imagem em escala de cinza, intensidade de pixel, faixa de 8 bits [0,255], vizinhanças e tratamento das bordas.}

\subsection{Operações algébricas e transformações de intensidade}
\placeholder{Descreva dissolve uniforme e não uniforme, negativo, limiarização, alargamento de contraste, transformação gamma e transformação logarítmica. Inclua as equações relevantes e explique seus parâmetros.}

\subsection{Histograma}
\placeholder{Explique histograma, expansão e equalização. Diferencie as duas operações e descreva o efeito esperado sobre o contraste.}

\subsection{Filtragem espacial e convolução}
\placeholder{Explique operações baseadas em vizinhança, convolução, máscaras, filtro da média e filtro da mediana. Diferencie filtros lineares e não lineares.}

\subsection{Detecção de bordas}
\placeholder{Apresente os operadores de Roberts, Sobel, Kirsch e Laplaciano. Explique as máscaras utilizadas, respostas direcionais quando aplicável, magnitude do gradiente e sensibilidade ao ruído.}

\subsection{Aguçamento}
\placeholder{Explique aguçamento baseado no Laplaciano e high boost. Discuta a combinação da resposta do operador com a imagem original e quando ocorre soma ou subtração.}

\subsection{Controle de contraste adaptativo}
\placeholder{Apresente a formulação utilizada na disciplina, explicando a constante c, a vizinhança n x n, a média local e o desvio padrão local.}

% ============================================================
% 3 MATERIAIS E MÉTODOS
% ============================================================
\section{Materiais e métodos}

\begin{orientacao}
Descreva \textbf{como} o sistema foi construído. Esta seção deve permitir que outra pessoa compreenda a estratégia de implementação. Informe ferramentas, organização do código, estruturas de dados, fluxo de processamento e decisões de projeto. Não apresente apenas capturas de tela.
\end{orientacao}

\subsection{Ambiente de desenvolvimento}
\placeholder{Informe linguagem, versão, IDE/editor, sistema operacional e bibliotecas utilizadas. Diferencie bibliotecas usadas para leitura/exibição daquelas que implementariam diretamente operações de processamento.}

\subsection{Arquitetura e organização do sistema}
\placeholder{Descreva os módulos, classes, funções ou componentes principais. Inclua, se necessário, um diagrama simples da arquitetura ou do fluxo de execução.}

\begin{figure}[H]
\centering
\fbox{\parbox[c][45mm][c]{0.85\textwidth}{\centering\color{muted}\textit{Insira aqui um diagrama da arquitetura ou fluxo do sistema, se aplicável.}}}
\caption{Arquitetura ou fluxo geral do sistema.}
\label{fig:arquitetura}
\end{figure}

\subsection{Implementação das operações}

\begin{orientacao}
Para cada grupo de operações, descreva a lógica implementada, os parâmetros de entrada, o processamento realizado e a saída. Destaque explicitamente o tratamento das bordas e dos valores fora de [0,255]. Evite colar grandes blocos de código; apresente apenas trechos essenciais quando ajudarem a explicar uma decisão.
\end{orientacao}

\subsubsection{Operações algébricas e transformações de intensidade}
\placeholder{Descreva como foram implementados dissolve, negativo, limiarização, alargamento de contraste, gamma e logaritmo.}

\subsubsection{Expansão e equalização de histograma}
\placeholder{Descreva como o histograma é calculado e como são realizadas a expansão e a equalização.}

\subsubsection{Filtros da média e mediana}
\placeholder{Descreva o percurso pela vizinhança, o tamanho parametrizável e o tratamento das bordas.}

\subsubsection{Convolução genérica}
\placeholder{Explique a implementação de g = f*h + offset, a leitura/entrada da máscara n x m, o posicionamento da máscara, o tratamento das bordas e o tratamento de resultados negativos ou superiores a 255.}

\subsubsection{Roberts, Sobel, Kirsch e Laplaciano}
\placeholder{Explique como cada operador foi implementado. Para operadores direcionais, descreva como as respostas das máscaras são combinadas para gerar a resposta final.}

\subsubsection{Aguçamento e high boost}
\placeholder{Descreva como a resposta do Laplaciano é combinada com a imagem original e como foi implementado o high boost.}

\subsubsection{Controle de contraste adaptativo}
\placeholder{Descreva a implementação da média e do desvio padrão locais, a parametrização de c e da vizinhança n x n.}

% ============================================================
% 4 RESULTADOS
% ============================================================
\section{Resultados}

\begin{orientacao}
Apresente evidências do funcionamento do sistema. Para cada funcionalidade relevante, mostre a imagem original, parâmetros utilizados e imagem resultante. As figuras devem possuir legenda e ser mencionadas no texto. Não basta inserir imagens: descreva o que foi observado.
\end{orientacao}

\subsection{Interface e fluxo de utilização}
\placeholder{Apresente a interface do sistema e explique como o usuário seleciona imagens, operações, máscaras e parâmetros.}

\begin{figure}[H]
\centering
\fbox{\parbox[c][55mm][c]{0.85\textwidth}{\centering\color{muted}\textit{Substitua por uma captura da interface do sistema.}}}
\caption{Interface principal do sistema desenvolvido.}
\label{fig:interface}
\end{figure}

\subsection{Transformações de intensidade e histogramas}
\placeholder{Apresente testes representativos de negativo, limiarização, alargamento, gamma, logaritmo, expansão e equalização. Informe os parâmetros utilizados e interprete os resultados.}

\subsection{Filtragem espacial}
\placeholder{Apresente resultados dos filtros da média e mediana. Compare seus efeitos, especialmente em relação à suavização e ao ruído.}

\subsection{Testes de convolução}

\begin{orientacao}
Inclua obrigatoriamente os testes de convolução definidos na avaliação: filtro de aguçamento com diferentes valores de c e d, filtros de relevo e máscaras de detecção de bordas. Para cada teste, informe a máscara, parâmetros, tratamento aplicado a valores fora de [0,255] e interpretação do resultado.
\end{orientacao}

\subsubsection{Filtro de aguçamento}
\placeholder{Apresente os resultados para diferentes valores de c e d, incluindo c=d=1. Relacione a soma dos coeficientes ao comportamento observado.}

\subsubsection{Filtros de relevo}
\placeholder{Compare as máscaras fornecidas e explique como a orientação dos coeficientes alterou a aparência do relevo.}

\subsubsection{Máscaras de detecção de bordas}
\placeholder{Apresente e interprete os resultados obtidos com as máscaras indicadas na avaliação.}

\subsection{Roberts, Sobel, Kirsch e Laplaciano}
\placeholder{Apresente resultados para uma mesma imagem sempre que possível. Mostre respostas intermediárias/direcionais e a resposta final. Compare visualmente os operadores.}

\subsection{Aguçamento pelo Laplaciano e high boost}
\placeholder{Apresente a resposta do Laplaciano separadamente e a imagem após combinação com a original. Inclua também um teste de high boost, informando os parâmetros.}

\subsection{Controle de contraste adaptativo}
\placeholder{Apresente testes com diferentes valores de c e/ou dimensões de vizinhança e explique o efeito dessas alterações.}

% ============================================================
% 5 DISCUSSÃO
% ============================================================
\section{Discussão}

\begin{orientacao}
Interprete criticamente os resultados. Compare métodos, explique diferenças observadas, discuta limitações e relacione o comportamento do sistema à teoria. Esta seção não deve repetir a seção de Resultados.
\end{orientacao}

\placeholder{Discuta, entre outros aspectos: média versus mediana; Roberts versus Sobel versus Kirsch versus Laplaciano; detecção de bordas versus aguçamento; expansão versus equalização; influência de parâmetros; efeitos do tratamento das bordas; problemas encontrados durante a implementação e como foram resolvidos.}

\subsection{Limitações e dificuldades encontradas}
\placeholder{Descreva limitações técnicas, casos em que os resultados não foram satisfatórios, erros encontrados durante o desenvolvimento e decisões tomadas para corrigi-los.}

% ============================================================
% 6 CONCLUSÃO
% ============================================================
\section{Conclusão}

\begin{orientacao}
Retome o objetivo do trabalho e sintetize os principais resultados alcançados. Indique o que foi efetivamente implementado, quais conhecimentos foram consolidados e quais limitações permaneceram. Não introduza resultados novos nesta seção.
\end{orientacao}

\placeholder{Escreva aqui a conclusão do trabalho.}

% ============================================================
% 7 IA
% ============================================================
\section{Declaração de uso de ferramentas de Inteligência Artificial}

\begin{orientacao}
Informe, de forma objetiva, se ferramentas de IA generativa foram utilizadas como apoio. Caso tenham sido utilizadas, descreva a finalidade (por exemplo: consulta conceitual, depuração, revisão textual ou auxílio na compreensão de mensagens de erro). A declaração não substitui a responsabilidade do aluno sobre o código, o texto e os resultados apresentados.
\end{orientacao}

\begin{campo}
\textbf{Ferramentas utilizadas:} \placeholder{Informe ``não utilizadas'' ou liste as ferramentas.}\\[3mm]
\textbf{Finalidade do uso:} \placeholder{Descreva brevemente como foram utilizadas.}\\[3mm]
\textbf{Observações:} \placeholder{Se necessário, detalhe quais partes receberam auxílio.}
\end{campo}

% ============================================================
% REFERÊNCIAS
% ============================================================
\section{Referências}

\begin{orientacao}
Liste as fontes efetivamente utilizadas na fundamentação e na implementação. Utilize um padrão bibliográfico consistente. Inclua livros, artigos, documentação técnica e materiais da disciplina quando citados no texto.
\end{orientacao}

\placeholder{Exemplo: SOBRENOME, Nome. Título da obra. Edição. Local: Editora, ano.}

% ============================================================
% APÊNDICE
% ============================================================
\appendix
\section{Informações complementares}

\begin{orientacao}
Se necessário, utilize o apêndice para tabelas extensas, máscaras adicionais, trechos curtos de código essenciais à compreensão, resultados complementares ou instruções de execução. Não reproduza integralmente o código-fonte; o código completo deverá permanecer no repositório indicado.
\end{orientacao}

\placeholder{Conteúdo opcional. Remova esta seção caso não seja necessária.}

\end{document}
